\chapter[Sermon 19. He Proceeds in Reciting Most Great Rewards.]{He Proceeds in Reciting Most Great Rewards.}
\label{sermon:019}
\markright{Sermon 19. He Proceeds in Reciting Most Great Rewards.}

He who overcomes, I will make a pillar in the temple of my God, and he shall go no more out. And I will write upon him the name of my God, and the name of the city of my God, New Jerusalem, which comes down out of Heaven from my God, and I will write upon him my new name. Let him who has ears hear what the Spirit says to the congregations.

Our Lord proceeds in recounting much more ample rewards, which he would give to those who overcome: and so he tempers his words, that we may easily perceive that this promise does not only pertain to the congregation of Philadelphia, but to all the churches in the whole world, yes, and to every of the faithful. And as we have oftentimes repeated already (for I am not ashamed to repeat, saying that the Lord himself so greatly urges the victory), again we say, that those things are promised not to such as fight lightly or negligently (for diverse fight, and by and by run away) but to those that overcome and persevere to the end. For our life is a warfare upon earth, which Job also has confessed. The soldier has a sure purpose to overcome his enemies. Our enemies are the devil, the world, and the flesh. Against these we must earnestly fight, of no other intent but that we may overcome. The Apostle in the 6th to the Ephesians described the armor of the faithful. In victory the saints consider sincerity and integrity, that we lose nothing of the known truth; but let us retain the pure word of God and sincere faith, and let us keep our bodies and our souls clean from all pollution, and that to our lives’ end. He propounds most ample rewards by promise. Whereby he alludes to the manner of Greeks and Romans, who decreed images to those who deserved well of the common wealth, in which also they wrote their virtues, for whose cause they were set up either in the court or marketplace or else where. They seemed by this means to deliver to their posterity as it were by hand the glory of their elders, which they made also as it were everlasting. Otherwise the use of this vocable Column, or pillar is diverse. Jeremiah is called by God a pillar for his constancy. The Apostles are called by Paul in the 2nd to the Galatians chief pillars, for their excellence, and that the churches leaned upon them, for the preaching of the truth. The church itself also is called the pillar and base of truth, for as much as it is grounded upon the sure rock Christ. In the Temple of Solomon stood two columns or pillars, figures of the everlasting kingdom of Christ, and of the holy church. In this place a pillar is taken for a man, in glory and beauty excellent. For he says not that he will erect a pillar for a godly man: but I will, says he, make him a pillar, that is to say, I will beautify him with honors and glory everlasting.

\index{Arethas of Caesarea}But where shall this pillar be established? where shall the glory of the saints be renowned? Not in a royal court or marketplace, but in the temple of my God. And the temple of God is heaven itself, and in this world, the holy Church. Therefore, he will be glorious in the church of the saints, both those engaged in warfare [militaunt] and those victorious [triumphaunt]. Although the faithful may hear evil in this world, this world will perish, Christ will reign forever, and the saints will reign with him. Their glory will remain forever and ever. And where he says, “of my God,” Arethas explains and says: This saying “of my God” does not diminish the divine nature seen in Christ, but establishes, as I may say, the consubstantiality. For it declares the union of two natures that are in the person of our Lord Jesus, namely his deity and humanity, even after understanding, yet not confusedly, to be inseparable. For they answer one another mutually, because of the assumption of the human nature, the divine to the human, and likewise the human to the divine properties. \&c.

Moreover, the lasting and steadfast nature of the glory of the saints and faithful is signified, where it is added, “and he shall go no more out.” For often pillars are broken and cast down, and the renown of great men perishes and fades away. But Christ promises to those who overcome that they shall never be cast out of the fellowship of saints, neither that the glory of the faithful should be obscured at any time. And thus far concerning the pillar itself.

\index{Rome}Here follows a description of the inscription on the pillar, noting what sort it will be. Three things chiefly are written on the Saints: the name of God, the name also of the city of God, and the new name of God or of Christ. Which we shall discuss in order. First, the name of God is ascribed to the godly, that is, they themselves are called by the name of God, and are the children and heirs of God. This is discussed at length in 1 John, chapters 1-8, and in Romans 8. And what can you devise to be more honorable than to be, and be called the son, not of a king nor emperor, but of the living God? But this same noble grace the Lord grants to those who overcome. This is found in 1 John 3 and 5. Secondly, to the overcomers is inscribed the name of the city of God: that is to say, the godly man is written in the number of the citizens of the city of God, and is truly a citizen of the city of God, I say even of the city of God. It was a great matter in times past to be a citizen of Rome, but it is far greater to be a citizen of the city of God. The citizens enjoy all privileges and commodities, finally the glory of the city. But this is greater and more than that it can be declared in few words.

But the church is the city of God: and the city of God is the church. This is presented here with three descriptions or titles, from which it is easy to judge what the church is, or what we should think of it. The church is the city of God. For just as a city is the fellowship of citizens, so the church is the communion of saints. The Prince of it is Christ, the head of the church. The entire world itself was a representation of this church, and the very setting up of the tents, in the midst of which was seen the Tabernacle, was a sign of the deity present, as if a cohabiter. \&c. For the Lord is in the midst of the church, as we have read in Leviticus 26 and 2 Corinthians 6. Secondly, the church is called New Jerusalem. For the old [?] was a representation of the new. This earthly church is New Jerusalem, that is to say, spiritual. Paul also affirms this in Galatians 4. For in the third place, this newness is explained. It is not built by humans, but comes down from heaven above. For unless we are born from above by spirit and by immortal seed, namely the word of God, we cannot be members of the church. And we are born by a spiritual regeneration, the children of Christ and of the church, as the Lord himself discusses at length in John 3.

1 Peter 3:1. And Paul, in his first letter to the Corinthians, chapters 3 and 4, will speak further of the New Jerusalem toward the end of this book. But understand this: what is the church of Christ, the fellowship of the faithful, regenerated by the word of God, and so forth.

Finally, in those who overcome, a new name is written, and indeed the new name of Christ. Not only that they should be called Christians of Christ, but because the name is a brief description of everything and nature, and a new name is promised, it follows that we should understand that people will be renewed, chiefly by glorifying. He promises therefore a glorifying to the godly. This is spoken elsewhere in Matthew 17:1, 1 Corinthians 15, Philippians 3, and 1 John 3. These most ample rewards the saints may truly expect if they fight to overcome.

Hereunto is annexed the customary affirmation, by which both this doctrine is applied and communicated to all churches throughout the world. And it is declared that it came not of men, as a vain thing, but of the very spirit of God, most true. This spirit the Lord grant us.
